%% ============================================================================
%% Fichier  : annexes/Ann_F_contacts.tex
%% Titre    : Contacts et Ressources
%% Version  : 1.0 -- novembre 2026
%% ============================================================================
%% Adresses de sites institutionnels données à titre indicatif, à vérifier
%% avant diffusion. Contact de NDAP SUKULU : celui fixé par l'auteur pour la
%% page de connexion du portail.
%% ============================================================================

\chapter{Contacts et Ressources}
\label{ann:contacts}

Les adresses ci-dessous sont celles des sites institutionnels à la date de
rédaction~; elles peuvent changer et sont à vérifier avant usage.

\section*{Autorités et institutions camerounaises et régionales}

\begin{tableaucours}{L{4.4cm}L{4.6cm}L{3.2cm}}{%
    \tableauHeader{Institution} &
    \tableauHeader{Rôle dans le cours} &
    \tableauHeader{Site}
}
ANTIC & Régulation de la sécurité électronique, audits de sécurité
(loi n°~2010/012) & \texttt{www.antic.cm} \\
MINPOSTEL & Destinataire des rapports d'audit de sécurité (art.~32) &
\texttt{www.minpostel.gov.cm} \\
COBAC (Secrétariat général) & Contrôle interne des établissements de
crédit de la CEMAC & \texttt{www.sgcobac.org} \\
OHADA & Actes uniformes (sociétés, droit comptable) &
\texttt{www.ohada.org} \\
Chambre des comptes de la Cour suprême & Contrôle juridictionnel des
comptes publics & · \\
Contrôle Supérieur de l'État & Contrôle administratif des services
publics & · \\
\end{tableaucours}

\section*{Organisations professionnelles et normatives}

\begin{tableaucours}{L{4.4cm}L{4.6cm}L{3.2cm}}{%
    \tableauHeader{Organisation} &
    \tableauHeader{Ressources} &
    \tableauHeader{Site}
}
IIA & Normes mondiales d'audit interne, modèle des trois lignes &
\texttt{www.theiia.org} \\
ISACA & COBIT, ITAF, certification CISA & \texttt{www.isaca.org} \\
COSO & Contrôle interne, gestion des risques, fraude &
\texttt{www.coso.org} \\
INTOSAI et IDI & ISSAI, GUID~5100, manuels de mise en œuvre &
\texttt{www.intosai.org}, \texttt{www.idi.no} \\
ISO & Normes ISO et ISO/IEC & \texttt{www.iso.org} \\
NIST & SP~800-53, 800-53A, AI RMF & \texttt{www.nist.gov} \\
UNESCO & Recommandation sur l'éthique de l'IA & \texttt{www.unesco.org} \\
\end{tableaucours}

\section*{Recherche}

Les articles A0 (\texttt{hal-05701007}) et A1 (\texttt{hal-05701270}) sont
déposés sur l'archive ouverte HAL (\texttt{hal.science}). Les travaux en
préparation y seront déposés à leur tour.

\section*{NDAP SUKULU}

Les fiches d'auto-évaluation et les dix labs du cours sont disponibles sur
NDAP SUKULU, sous le code MSN-2S3-MSI05. NDAP SUKULU est une propriété
de MaletYon, mise à la disposition des étudiants et des praticiens
souhaitant se perfectionner~: \texttt{maletyon@proton.me}.
